\documentclass[11pt]{article}
% Shared preamble for all daily exercises
% Update this file to change settings (padding, parindent, etc.) for all exercises

\usepackage[margin=0.75in]{geometry}
\usepackage{amsmath}
\usepackage{amssymb}

% Custom settings
\setlength{\parindent}{0pt}
\setlength{\parskip}{0.2cm}

\title{Daily Exercise}
\author{Dhruman Gupta}
\date{21st March}

\begin{document}

% Common header for daily exercises
% This file can be included in other LaTeX documents

\maketitle

\noindent
\textbf{Course:} CS-3330, Theory of Computation

\vspace{0.2cm}

\noindent
\textbf{Question:}\noindent 
Show that recursively enumerable languages are closed under quotienting.

\textbf{Answer:}

Take right quotient to mean
\[
L_1 / L_2 = \{x \mid \exists y \in L_2 \text{ such that } xy \in L_1\}.
\]

Let $A$ and $B$ be recursively enumerable languages. We will show that $A/B$ is recursively enumerable.

Since $A$ and $B$ are recursively enumerable, there exist Turing machines $M_A$ and $M_B$ that recognise $A$ and $B$ respectively.

That is:
\begin{itemize}
    \item if $w \in A$, then $M_A$ accepts $w$ in finite time,
    \item if $w \notin A$, then $M_A$ may either reject or loop forever,
    \item if $w \in B$, then $M_B$ accepts $w$ in finite time,
    \item if $w \notin B$, then $M_B$ may either reject or loop forever.
\end{itemize}

We construct a nondeterministic Turing machine $N$ that recognises $A/B$ as follows.

On input $x$:
\begin{itemize}
    \item Nondeterministically guess a string $y$.
    \item Run $M_B$ on input $y$.
    \item Run $M_A$ on input $xy$.
    \item If both machines accept, then accept.
\end{itemize}

Now, if $x \in A/B$, then by definition there exists some $y \in B$ such that $xy \in A$. On the branch where $N$ guesses this $y$, the machine $M_B$ accepts $y$ and the machine $M_A$ accepts $xy$. Hence $N$ accepts $x$.

On the other hand, if $x \notin A/B$, then for every string $y$, at least one of the following holds:
\begin{itemize}
    \item $y \notin B$, or
    \item $xy \notin A$.
\end{itemize}

So on every branch, at least one of $M_B$ and $M_A$ does not accept. Hence no branch accepts, and so $N$ does not accept $x$.

Thus $N$ recognises $A/B$, so $A/B$ is recursively enumerable.

Therefore recursively enumerable languages are closed under quotienting.

\end{document}
